\section{The class sum on the prescribed higher cover}

% BEGIN T73C1
\paragraph{T73C1. The value of a short series on a framed point.}
Work on the upper level of one pinned tuple, with its given coefficient field,
connection, lattice $\Lambda$, odd auxiliary integer $N=N_c$, polynomial
$P$ in the indicated harmonic highest-weight space and residue function $w$.
Write $r=p+q$, $\mathcal Q=(\mathrm{LevQ})^n$ and
$d_\Lambda=[\Lambda\cap\mathcal O_E^n:N\mathcal O_E^n]^r$.
The residue set is finite since the half-level ideal has finite index.
The auxiliary choices are an enumeration $z_i$ of the source's isotropic
torsion subgroup, $N$-division points $t_{q,b}$ of the marked level
points in the specified half-level primary torsion, and the translated
sections $F_{qib}$ of the $N^2$-th power of the given theta line.
For the original class-sum normalization $\beta=1$, the claimed value is
\[
 \frac{1}{d_\Lambda\,|\mathcal Q|\,N^{\deg P}}
 \sum_{q\in\mathcal Q} w(q)
 \sum_i\sum_{m\in\operatorname{supp} P} P_m
       \prod_{b=1}^n \nabla_{m_b}F_{qib}.
\]
Here $m_b$ is the column multi-index and all coefficients map to the
actual test algebra. These choices can be made after a faithfully
flat extension. Indeed multiplication by $N$ on finite torsion is
surjective onto the required marked torsion, since $N$ is invertible
over the coefficient field. Split the primary components to choose
the prescribed half. Trivialize the finite étale isotropic group on
a common cover. Its geometric description comes from applying the
relations defining $\mathrm{InZh}$ to the marked Tate module:
the coefficients of $N\Lambda^\vee$ annihilate the tuple modulo
$N^2$, leaving the dual quotient of order
$[\Lambda:N\mathcal O_E^n]^r$ since $\Lambda^\vee=\Lambda$.
Thus a geometric injection indexed by this cardinality enumerates
the whole group. Both assertions may be checked on geometric fibers
of the finite étale schemes. The rigidified symmetric theta group
supplies the prescribed translations of the section. For odd order
its symmetric lift is unique: the two opposite lifts are related by
inversion and the square map is bijective on the odd roots of unity.
This is the odd-torsion theta-group construction of Mumford,
\emph{On the equations defining abelian varieties I}, \S1.

Here are the two choices that can affect this formula. First, in
logarithmic coordinates for the fixed tangent eigenframe, compare
two rigidified formal theta frames. Their residual biextension ratio
has logarithm bilinear, by additivity in both variables, with
coefficients supported on conjugate eigenspaces by the adjoint
identity. The constant is fixed by rigidification. In a column the
corrected germs therefore differ by $\exp(N^2 Q)$ for a quadratic
form $Q$ pairing conjugate directions. This follows as well by
subtracting the two normalized cocycle identities of CNC6--CNC12
in the same logarithmic coordinates. In the finite jet pairing
multiplication by $\sum_b Q(X_b)$ is adjoint to the corresponding
linear combination of the mixed Laplacians
$\sum_b\partial_{u,b}\partial_{v,b}$ on $P$. Each kills $P$.
Expanding the exponential through degree $\deg P$ shows that only
the constant term contributes. Thus the contracted value, including
its sum over monomials, commutes with arbitrary coefficient base
change. Individual uncontracted derivatives need not have this
property.

Second, on a common local base write $t'_{q,b}=t_{q,b}+u_b$.
The $u_b$ are $N$-torsion in the same half. The tuple $u$ belongs
to the isotropic group: at primes supporting the half-level,
$\Lambda$ is the standard lattice by $\mathrm{TrivAt}$ and
$Nu=0$, and the other components of $u$ vanish. Translation
$z'_i\mapsto z'_i+u$ permutes the first enumeration. The symmetric
theta lift in column $b$ changes by a scalar $r_b$ whose square is
the Weil commutator $e_{N^2}(u_b,z'_{ib})$ (in the inverse convention
both factors invert) and whose $N^2$-th power is one. Indeed moving
the half-torsion lift through the translation has precisely that
commutator after squaring. The commutator with $t_{q,b}$ itself
is trivial by the half-level isotropy. The product of the displayed
Weil commutators is one, by isotropy of the tuple group for the
product pairing. In lattice terms this is the integrality of the
Hermitian pairing on the self-dual $\Lambda$ and the adjointness
of the two level halves. Thus $\prod_b r_b$ has square and
$N^2$-th power one. As $N$ is odd it equals one. Connection values
are linear in the section, so the product in each monomial is
unchanged after reindexing. This argument on local rings gives
equality everywhere by localization at all maximal ideals.

Apply the two comparisons to the two pullbacks of the chosen
formula to the tensor square of the faithfully flat test algebra.
The faithfully flat equalizer for functions descends it to a
unique element in the original algebra. Comparing two choices
on their tensor product proves that this element fulfills the
literal $\mathrm{ShortValAt}$ predicate uniquely and that it is
natural for every framed-point base change. This applies in
particular to every algebraically closed field used by
$\mathrm{IsThClsD}$.
% END T73C1

% BEGIN T73C2
\paragraph{T73C2. Cancellation of the two level characters and descent.}
Fix a pinned datum $O$ and tuple $t$. Put $M=O.c'.F$ and
$d=O.d_\delta(t)$. Use the prescribed isomorphism $e_F$ of coefficient
fields for transport to the upper level. On the quotient lattice
$\mathrm{LaI}$, take the determinant of an automorphism, reduce at
each $P$ in $R.Q(t)$ and multiply the prescribed residue characters
$R.\xi_R(t)_P$. This defines $\psi$ over the field $M$: each factor
occurs in the pinned field by RTG2--RTG4, and every such $P$ contains
the half-level ideal. The determinant on a finite projective module
of constant rank is intrinsic, equivalently its determinant in a
basis over the finite semilocal quotient ring.

For a move between two upper markings of the same lower point,
the finite residue function of the tuple transforms by this
determinant character on the characteristic variable. The factors
at good places are fixed by the lower-level congruence; at a type
place the determinant of an integral matrix lifting the move reduces
to the determinant of its reduced matrix. The multiplier on the
Fourier-transformed function $w$ with the inverse change of summation
variable is $\psi^{-1}$. Indeed the finite perfect pairing identifies
the action on its two arguments by inverse adjoints, so changing the
residue variable permutes the sum and leaves its additive kernel
unchanged. These are the source's $\mathrm{DatEq}$ equations at
each of the chosen representatives, with the $w$ of RTG11.

The cover section $c_\delta$ has the opposite multiplier $\psi$.
At the distinguished complex realization the pinned value is
$c_0\,\xi(\det h)$ on the marked frame. Two lifts differ by the
level move just considered and a rational unitary element. The latter
has character one, and on the former the determinant character
has precisely the finite residue product defining $\psi$.
This gives the ratio $\psi$ on the corresponding components.
For clarity concerning other geometric points, some positive power
of $c_0$ and of every $\xi(\det h)$ is one (use the product of their
orders). The marked dictionary reaches the underlying point of every complex frame:
use the period cover and the finite and real isometries of that marked
lattice relative to the specified frame. A weight-zero section is
independent of the frames over that point (the theta scaling has
weight zero and a regular function on the complete flag fiber is constant).
Its power is thus one everywhere. Holomorphy and connectedness of the
period domain now force constancy on each component. Comparing the
markings $h$ and $hu$ for any integral lift $u$ gives the stated ratio.
The resulting character identity is an equality of regular
functions on the level-move scheme. It descends from the chosen
complex embedding by faithful flatness and extends to every
algebraically closed field. The same power identity gives modulus
one at every good embedding. These are the $\mathrm{PsiTie}$ and
$\mathrm{UnitC}$ properties needed below.

For an upper lift $y_2$ of a geometric lower frame $y$ multiply
the unique short value of T73C1 by $\operatorname{ev}(c_\delta,y_2)$.
Changing the marking multiplies the two factors by $\psi^{-1}$
and $\psi$. The comparison holds also between a lift made with
the canonical level-forgetting map and a lift for the stored
map in $d$: both normalized identifications of the polarized
abelian variety respect the marked lower level, and their
composite is one of the integral upper marking moves in
$\mathrm{DatEq}$ and $\mathrm{PsiTie}$. At $A_2=A$ the extra
$A$-marking agrees as well. These identities therefore compare
the two products without identifying the two maps.

For a reduced original test algebra $S$, use the canonical
finite étale surjective level map for the specified pair
of levels to choose an algebra $S'$ on which the
underlying point lifts. Its geometric fibers are the allowed
upper markings: they lift because the torsion orders are
invertible and the level congruence map is surjective on
compatible markings. Carry the frame by the polarization and
universal-family isomorphism over this canonical map.
The only coherent section base change used here and in the
framed-point return of TDR13 concerns the first theta power.
For the polarized abelian scheme in characteristic zero, the
theta line on each geometric fiber is ample with trivial
canonical bundle, hence has no higher cohomology by Kodaira
vanishing. Riemann--Roch and principal polarization give
dimension one in degree zero. The proper smooth pushforward
of this line is perfect and commutes with derived base change;
on local affine charts a minimal finite free complex is
concentrated in degree zero by these fiber dimensions. Thus
the pushforward is a line and every base change, including
nonflat ones, pulls this line back with its generator. The
tangent eigenframes pull back by the Lie bundles. These facts
give precisely the framed-point comparisons used here.
The product defines an element of $S'$. Its two images in
$S'\otimes_S S'$ agree at every geometric point: choose
a lift for the original $d$ by $\mathrm{IsLev}$ and compare
each canonical pullback to it. The tensor product is reduced
because it is étale over reduced $S$. The faithfully flat
equalizer yields one element of $S$. The same
geometric criterion proves naturality in $S$ and the Borel character
law: the lower frame move lifts to the upper frame, the harmonic
highest-weight polynomial transforms by the prescribed row character,
the product of the $n$ theta columns gives the prescribed section
weight, and $c_\delta$ has weight zero. Thus the descended rule
belongs to $\mathrm{sec}$ at weight $(a,n)$. TDR13 supplies the
framed-point return on arbitrary geometric fields, so the descent
identity holds at every lifted pair in the native definition.
% END T73C2

% BEGIN T73C3
\paragraph{T73C3. The unique section with the original class-sum values.}
For each actual class $j$ use T73C2 with $P_j=O.io(x_j)$ transported
by $e_F$ and with $w_j,\Lambda_j,N_j$ from this tuple. The pinned
range of $O.io$ is in the span of the harmonic highest-weight
polynomials. Their conditions at the specified weight are linear,
so each $P_j$ has those conditions, also after the coefficient
transport. Denote the one-class section by $f_j$. In the original
lower field $M$ set
\[
       f_x=\sum_j \operatorname{vol}\,\gamma_j^{-1}
                         \operatorname{ep}_j\, f_j .
\]
The class set is finite and each $\gamma_j$ is positive by RTG7.
This section has weight $(a,n)$. On every algebraically closed
extension with the two compatible coefficient maps, and on
\emph{every} lifted pair $y_2/y$, its value is
\[
 \operatorname{ev}(c_\delta,y_2)\,
      \sum_j\operatorname{vol}\,\gamma_j^{-1}\operatorname{ep}_j
            \,v_j(y_2).
\]
These are exactly the upper $\mathrm{thLaValAt}$ values with
$\beta=1$. This gives $\mathrm{IsThClsD}$ by its defining
formula. Conversely let any other section of this weight satisfy
that formula. Over every algebraically closed coefficient field
choose one permitted lift of each lower frame by $\mathrm{IsLev}$.
For each class, T73C1 identifies the two purported values at
this lift. Hence the two sections agree on every geometric frame.
For a reduced original test algebra $S$ and each prime
$\mathfrak p$, specialize the frame to an algebraic closure of
$\kappa(\mathfrak p)$ with the upper coefficient map chosen
by $e_F^{-1}$. Naturality makes the difference of the two values
belong to $\mathfrak p$ for every $\mathfrak p$. Reducedness
makes the difference zero. Thus the two natural rules agree
on all original tests. This uniqueness holds among all sections
of the weight in question.
% END T73C3

% BEGIN T73C4
\paragraph{T73C4. The class section is cuspidal.}
On each complex component at each good embedding use the principal
self-dual lattice in the supplied uniformization and its given
symplectic marking. Write the generator of the theta line as
$\vartheta$ in the marked period chart. Its $n$-fold external
product is a scalar multiple $\kappa$ of the tensor of the
level-one vacuum functions: these sections obey the same
lattice translation law and the space of sections of a principal
theta line is one dimensional. The scalar is nonzero and has
fixed absolute value on unit frames. Indeed the norm of the
given generator is fixed there by the unit-frame normalization,
and the squared norm of the vacuum, integrated on a period
parallelepiped, is the Gaussian integral after orthogonality
of the Fourier characters. Its period determinant is the
unit-frame normalizing factor (NFL5--NFL9).
A compatible upper marking can be chosen on the whole component.
Choose a lift over one point; the discrete finite lattice
marking persists as the period varies over the marked period
domain, and forgetting it recovers the same lower marking.
This selects the reached upper component also when the two
signatures are equal. There the weight-zero section $c_\delta$
has the constant value determined in T73C2.

Compute T73C1 on these analytic tori using
the same finite sums. Lift the marked level points to
characteristics $Y_q$ and the isotropic enumeration to $u_i$.
The two theta-group translations give the Heisenberg
operators for $Y_q$ and $u_i$ on the vacuum. In logarithmic
coordinates the harmonic contraction equals the derivative with
polynomial $P_j$ of the external product at zero: the extra
quadratic exponential disappears by the mixed Laplacian
calculation of T73C1. Multiplication by $N$ scales the
derivative by $N^{\deg P_j}$, canceled by the short-value
denominator since fixed row degrees make $P_j$ homogeneous.
The average over the finite isotropic subgroup is its projector,
so the sum over the enumeration contributes
$[\Lambda_j:N\mathcal O_E^n]^r$. The resulting finite Schwartz
function is the sum over $q$ of $w_j(q)$ times the Heisenberg
translate by $Y_q$ of this projected vacuum. Its scalar coefficient
is $[\Lambda_j:N\mathcal O_E^n]^r/(d_{\Lambda_j}|\mathcal Q|)$,
besides its class weight from T73C3. Thus at unit frames the
reading of $f_x$ is $\kappa$ times a finite linear combination
of ordinary theta series with these finite characteristics and
row-highest polynomials, including the fixed cover-section
factor. The tangent directions are the rows of the marked
period map. Normal convergence of the polynomial-weighted
Gaussian justifies differentiation.

The holomorphic reading and this combination have the same
Borel weight. Every frame is upper-triangular equivalent to
a unit frame by orthogonalization. Their absolute-value ratio
is therefore constant wherever the combination is nonzero.
There the ratio is holomorphic and locally constant by the
open mapping theorem. The identity principle extends one such
equality to the connected domain. If the combination vanishes
identically, so does the reading.

At a rational isotropic line choose an adapted symplectic
basis. Under the change of basis the finite characteristic
functions become finite sums of characteristics by Fourier
transformation. In the theta expansion the frequency along
the isotropic flow is a positive quadratic norm of the
coordinate in that direction. Zero frequency forces that
coordinate to vanish. At an embedding where the last row
weight is positive, every row-highest polynomial lies in the
span of products containing a full-height minor of that block:
the last fundamental-weight multiplicity is its last row
weight. On a zero-frequency summand the rows have a common
linear dependence given by the isotropic direction; in adapted
coordinates the linear combination of rows for that direction is
zero. Every such full-height minor vanishes. Absolute summability
allows averaging along the rational flow, which retains only
these zero terms. The finite characteristic theta functions
have moderate growth by the lattice Gaussian estimate and
its polynomial derivatives in Siegel coordinates. In the
definite signature there is no rational isotropic line. This
checks the cuspidal condition at every good embedding, using
the last-row positivity or definite alternative of the tuple
weights.
% END T73C4

% BEGIN T73C5
\paragraph{T73C5. Class-sum existence and uniqueness for every pinned record.}
Fix $\mathrm{CoreJ},R,\kappa,v_0$ with the original
$\mathrm{FWTie}$, level and family laws and prime agreement
of $\mathrm{clsEx\_of\_core}$. Take any $i$, any tuple-index
set and tuple function, and any $O$ containing the specified
embedded field with the exact $\mathrm{Pins}$. For each tuple
and each placed coefficient vector take the section of T73C3.
The cover data give $\mathrm{IsLev}$ and the given canonical
connection, while $\mathrm{RepsOK}$ gives the three lattice
conditions of T73C1. The marked lift and finite residue
polynomial give the inverse-character pair used in T73C2.
The base has $A=1$, the cover has $A_2=1$, and the uniformizations
after the prescribed coefficient extension come by composing
each good embedding with the base inclusion. The weight
agreement in $\mathrm{FWTie}$ gives $n=\operatorname{rank} V_i$
and the dominant row data. The source's $\mathrm{Dom}(a,n)$
includes that $n$ is even: the given $\kappa_i$ has even index,
and its index equals this rank by the same weight agreement.
This supplies the parity for the tensor-vacuum and theta
transformation formulas of T73C4. In the nondefinite case its last
positive label is the last entry of the positive row in the
original lift weight. Thus the last-row test in T73C4 applies
to this very $\kappa_i$. The cusp section in T73C4 and its
uniqueness in T73C3 give precisely $\mathrm{ClsEx}$.

If a permitted phase correction multiplies $c_\delta$ by the
transport of $u\in M^\times$, the equation of T73C3 multiplies
$f_x$ by $u$. All the short-value data are fixed. When the
corrected record satisfies the pins the same construction
applies, with every lifted pair and the original target space.
This verifies both occurrences of class-sum existence and
uniqueness around the phase correction in RTG16--RTG18.
The common pointwise comparison with the adelic lift and the
base $\mathrm{ClassSumsSpan}$ are additional assertions.
% END T73C5
